\documentclass[14pt,a4paper,onesize,oneside]{scrreprt}

\usepackage{fontspec}
\setmainfont{Liberation Serif}
\setsansfont{Liberation Sans}
\setmonofont{Liberation Mono}
\newfontfamily\cyrillicfont{Liberation Serif}

\usepackage{polyglossia}
\setdefaultlanguage{russian}

\usepackage{cmap}
\usepackage{microtype}
\usepackage{setspace}
\onehalfspacing

\usepackage[
  left=30mm,
  right=15mm,
  top=20mm,
  bottom=20mm
]{geometry}

\usepackage{booktabs}
\usepackage{hyperref}
\hypersetup{colorlinks=true,linkcolor=black,urlcolor=blue}

\setlength{\emergencystretch}{3em}
\setlength{\parindent}{1.25cm}
\newcommand{\pth}[1]{\path{#1}}
\newcommand{\weblink}[1]{{\urlstyle{same}\href{#1}{\nolinkurl{#1}}}}

\pagestyle{plain}

\begin{document}

\renewcommand*{\thesection}{\arabic{section}}
\setcounter{section}{0}

\begin{center}
  \vspace*{1.5cm}

  {\sffamily\large АННОТАЦИЯ\par}

  \vspace{1.2cm}

  {\LARGE\bfseries Предрелизная подготовка\\в системе MarLine\par}

  \vspace{1.8cm}

  \begin{tabular}{@{}p{6cm}@{\hspace{0.4cm}}p{9.6cm}@{}}
    Автор:                & Овечкин Андрей \\
    Научный руководитель: & Гориховский Вячеслав Игоревич \\
    Консультант:          & Дмитриевцев Алексей Сергеевич \\
    Репозиторий:          & \weblink{https://github.com/Iseyne/marline-reports} \\
    Дата:                 & 4 октября 2026 \\
  \end{tabular}

  \vspace{0.8cm}

\end{center}

\thispagestyle{empty}

\section{Аннотация}

MarLine — система дедупликации, построенная на границах, зависящих от
содержимого. Данные нарезаются на чанки, для каждого считаются признаки
подобия, похожие фрагменты собираются в кластеры и кодируются как дельты
относительно базового. Разработка идёт силами учебной команды и
охватывает шесть крейтов: построение эскизов, индекс с эвристическим
отбором кандидатов, кластеризацию, кодировщики дельт, хранилище на
\pth{chunkfs} и отдельный Palantir-конвейер.

Основная проблема сейчас даже не в алгоритмах. Модульных тестов
достаточно много, 281 штука, и подавляющее большинство приходится на
кодировщики и индекс, но end-to-end тестов нет вообще, а у бинарника
\pth{marline} тестов нет вовсе. Интеграционных файлов три, и полный
конвейер они не покрывают. Получается, что каждая функция проверена
по отдельности, но то, что все стадии правильно собираются в один прогон
по реальному файлу, не подтверждено ничем.

Попутно вылезло несколько вещей, из-за которых релиз в текущем виде
невозможен. Коэффициенты хеша признаков берутся из генератора случайных
чисел, поэтому два одинаковых запуска дают разные хеши и результаты
бенчмарков между собой не сравниваются. Конфигурационная сетка
в \pth{palantir_metrics} настроена, но фактически не применяется: все
варианты сводятся к одному, а метрика дедупликации задана константой
независимо от данных. Пример быстрого старта в README не компилируется,
и никто этого не замечает, потому что README не подключён к документации
и не проверяется \pth{doc}-тестами. Ни в одном крейте не заполнены поля,
которые нужны \pth{cargo} для публикации.

Цель работы — довести проект до состояния, в котором его можно
нормально выпустить.

Результатом планируется набор интеграционных и end-to-end тестов,
закрывающих полный цикл от нарезки до дедупликации, исправление
найденных дефектов и приведение документации в порядок — с рабочими
примерами вместо тех, что не компилируются. Отдельно стоит вопрос
с зависимостями: внешний \pth{chunkfs} подключён по адресу \pth{git} без
фиксации ревизии, поэтому ее версия может смениться без нашего участия,
и сборка станет невоспроизводимой сама по себе.

\section{Что уже проверено}

Обследование кода показало следующее.

\begin{center}
\begin{tabular}{@{}lrrr@{}}
  \toprule
  \textbf{Крейт} & \textbf{Модульные} & \textbf{Интеграционные} & \textbf{End-to-end} \\
  \midrule
  \pth{marline_index}    & 80  & 0 & 0 \\
  \pth{marline_scrub}    & 165 & 7 & 0 \\
  \pth{marline_palantir} & 19  & 2 & 0 \\
  \pth{marline_sketcher} & 8   & 0 & 0 \\
  \pth{marline}          & \textbf{0} & 0 & 0 \\
  \pth{marline_delta}    & 0   & 0 & 0 \\
  \midrule
  \textbf{Итого}         & \textbf{281} & \textbf{9} & \textbf{0} \\
  \bottomrule
\end{tabular}
\end{center}

Модульное тестирование для алгоритмической части сделано неплохо,
и трогать его не нужно. Проблема в другом: каталог \pth{marline_delta}
пустой, непонятно, зачем он вообще нужен, а поведение собранного
бинарника не проверяет ни один тест.

Сценарий непрерывной интеграции тоже стоит расширить. Сейчас он
срабатывает только на \pth{push}, из-за чего ветку с pull request можно
слить, не дождавшись проверок, и в нём нет ни \pth{cargo publish
--dry-run}, ни сборки документации, которая бы поймала неработающие
ссылки.

\section{Направления работы}

\begin{itemize}
  \item Сделать поведение системы проверяемым: тесты на полный
        цикл работы конвейера и на восстановление данных без потерь.
  \item Убрать недетерминизм, иначе любые замеры теряют смысл.
  \item Привести в порядок документацию так, чтобы примеры
        собирались и проверялись автоматически.
  \item Почистить репозиторий и настроить CI так, чтобы сбой
        проверок останавливал выпуск релиза.
\end{itemize}

\end{document}
